% chapters/requirements/rf_exams.tex
% RF-6: Exámenes y modelos
% ============================================================================

\item \textbf{Exámenes y modelos.} \label{rf:exams}
Un examen es la entidad que representa una prueba evaluable dentro de una asignatura (por
ejemplo, <<Parcial~1 de Álgebra>>). Cada examen tiene una lista de alumnos convocados, un
conjunto de reglas de asignación de correctores y al menos un modelo. Los modelos definen
variantes estructurales del examen: número de problemas, puntuaciones, zonas de
reconocimiento y un \acs{pdf} en blanco de referencia. Para cada alumno convocado existe
una única instancia de examen, independientemente del modelo. Cuando los modelos comparten
estructura, el sistema permite copiar la configuración de un modelo como plantilla base.

\begin{functional}

    % RF-6.1 ----------------------------------------------------------------
    \item \textbf{Creación de examen.} \label{rf:exams:create}
    El sistema debe permitir al coordinador o a un profesor con el permiso
    \texttt{can\_create\_exams} crear un examen dentro de una asignatura.

    \textbf{Entrada:} petición POST al recurso de exámenes de la asignatura con el nombre
    del examen.

    \textbf{Proceso:} el sistema crea el examen asociado a la asignatura (que pertenece al
    curso activo de la organización). Se crea automáticamente un primer modelo con
    identificador <<A>> y sin problemas definidos. La lista de convocados se inicializa
    vacía.

    \textbf{Salida:} respuesta con los datos del examen creado, incluyendo el modelo inicial,
    y código \acs{http}~201.

    % RF-6.2 ----------------------------------------------------------------
    \item \textbf{Modificación de examen.} \label{rf:exams:update}
    El sistema debe permitir modificar los datos de un examen existente.

    \textbf{Entrada:} petición PATCH al recurso del examen con los campos a modificar
    (nombre y permisos del alumno sobre la instancia).

    \textbf{Proceso:} el sistema actualiza los datos del examen. Solo se permite la
    modificación si el examen no tiene instancias en estado publicado o posterior. Si se
    modifican los permisos del alumno, el cambio se aplica a todas las instancias del
    examen.

    \textbf{Salida:} respuesta con los datos actualizados del examen.

    % RF-6.3 ----------------------------------------------------------------
    \item \textbf{Eliminación de examen.} \label{rf:exams:delete}
    El sistema debe permitir eliminar un examen que no tenga instancias asociadas.

    \textbf{Entrada:} petición DELETE al recurso del examen.

    \textbf{Proceso:} el sistema verifica que el examen no tiene instancias creadas. Si las
    tiene, la eliminación se rechaza. Si no las tiene, se eliminan los modelos, zonas de
    reconocimiento, problemas y rúbricas asociados, y se eliminan los ficheros \acs{pdf} en
    blanco de MinIO.

    \textbf{Salida:} respuesta con código \acs{http}~204 si se elimina correctamente, o
    código \acs{http}~409 si tiene instancias.

    % RF-6.4 ----------------------------------------------------------------
    \item \textbf{Gestión de la lista de convocados.} \label{rf:exams:call_list}
    El sistema debe permitir al coordinador o a un profesor con el permiso
    \texttt{can\_manage\_call\_list} definir la lista de alumnos convocados a un examen.

    \textbf{Entrada:} petición PUT al recurso de convocados del examen con la lista de
    identificadores de alumnos que deben estar convocados.

    \textbf{Proceso:} el sistema recibe la lista completa de identificadores de alumnos
    convocados. Valida que todos los identificadores correspondan a estudiantes activos de
    la asignatura con \texttt{SubjectMembership} activa. Reemplaza la lista de convocados
    anterior por la nueva. Solo pueden convocarse alumnos que tengan grupo asignado.

    \textbf{Salida:} respuesta con la lista actualizada de convocados y el número total.

    % RF-6.5 ----------------------------------------------------------------
    \item \textbf{Creación de modelo con copia de estructura base.}
    \label{rf:exams:create_model}
    El sistema debe permitir crear un modelo dentro de un examen, copiando opcionalmente la
    estructura de otro modelo existente.

    \textbf{Entrada:} petición POST al recurso de modelos del examen con un identificador de
    modelo (cadena alfanumérica, por ejemplo <<B>>) y, opcionalmente, el identificador de
    otro modelo del que copiar la estructura.

    \textbf{Proceso:} si se indica un modelo fuente, el sistema copia su lista de problemas
    con sus puntuaciones, rúbricas, zonas de reconocimiento y referencia al \acs{pdf} en
    blanco. Si no se indica fuente, el modelo se crea vacío.

    \textbf{Salida:} respuesta con los datos del modelo creado, incluyendo los problemas
    copiados si procede, y código \acs{http}~201.

    % RF-6.6 ----------------------------------------------------------------
    \item \textbf{Modificación de modelo.} \label{rf:exams:update_model}
    El sistema debe permitir modificar la estructura de un modelo: sus problemas,
    puntuaciones y zonas de reconocimiento.

    \textbf{Entrada:} petición PATCH al recurso del modelo con los datos a modificar.

    \textbf{Proceso:} el sistema actualiza la estructura del modelo.

    \textbf{Salida:} respuesta con los datos actualizados del modelo.

    % RF-6.7 ----------------------------------------------------------------
    \item \textbf{Carga de \acs{pdf} en blanco.} \label{rf:exams:blank_pdf}
    El sistema debe permitir cargar el \acs{pdf} en blanco de referencia de un modelo,
    necesario para definir las zonas de reconocimiento.

    \textbf{Entrada:} petición POST o PUT al recurso de \acs{pdf} en blanco del modelo con
    el fichero \acs{pdf}.

    \textbf{Proceso:} el sistema valida que el fichero es un \acs{pdf} válido y no excede el
    tamaño máximo configurado por organización. El fichero se almacena en MinIO en un
    \dfn{bucket} dedicado a plantillas de examen. Si el modelo ya tenía un \acs{pdf} en
    blanco, se sobreescribe. Se permite reutilizar el mismo \acs{pdf} entre modelos del
    mismo examen mediante una referencia.

    \textbf{Salida:} respuesta con la confirmación de la carga y la referencia al fichero
    almacenado.

    % RF-6.8 ----------------------------------------------------------------
    \item \textbf{Definición de zonas de reconocimiento.} \label{rf:exams:zones}
    El sistema debe permitir definir regiones rectangulares sobre el \acs{pdf} en blanco de
    un modelo, asociadas a atributos del examen, usuario, asignatura o curso.

    \textbf{Entrada:} petición POST al recurso de zonas del modelo con los datos de la
    zona: coordenadas del rectángulo ($x$, $y$, ancho, alto), número de página y atributo
    asociado (nombre del alumno, \acs{dni}, \acs{nia}, identificador de modelo,
    identificador de examen, problema, código de asignatura, etc.).

    \textbf{Proceso:} el sistema valida que las coordenadas estén dentro de los límites de la
    página indicada del \acs{pdf} en blanco y que el atributo asociado sea válido. Se
    almacena la zona vinculada al modelo.

    \textbf{Salida:} respuesta con los datos de la zona creada y código \acs{http}~201.

    % RF-6.9 ----------------------------------------------------------------
    \item \textbf{Definición de problemas.} \label{rf:exams:problems}
    El sistema debe permitir definir los problemas o partes de un modelo, con su
    puntuación máxima y, opcionalmente, zonas de reconocimiento asociadas que pueden
    abarcar varias páginas.

    \textbf{Entrada:} petición POST al recurso de problemas del modelo con los datos: nombre
    o número del problema (por ejemplo, <<Problema~1>>), puntuación máxima (decimal) y,
    opcionalmente, una lista de zonas (cada una con su página y coordenadas).

    \textbf{Proceso:} el sistema crea el problema y, si se proporcionan, sus zonas
    asociadas. Un problema puede tener cero zonas (si la respuesta es de longitud
    indefinida) o varias (si se distribuye en múltiples regiones o páginas). Las
    puntuaciones pueden ser negativas (penalizaciones).

    \textbf{Salida:} respuesta con los datos del problema creado, sus zonas y código
    \acs{http}~201.

    % RF-6.10 ---------------------------------------------------------------
    \item \textbf{Configuración de permisos del alumno sobre el examen.}
    \label{rf:exams:student_perms}
    El sistema debe permitir configurar qué información puede ver el alumno sobre su
    instancia una vez publicada.

    \textbf{Entrada:} petición PATCH al recurso del examen con los permisos del alumno a
    modificar.

    \textbf{Proceso:} el sistema actualiza los permisos configurables: ver \acs{pdf},
    ver \acs{pdf} extendido (incluso tras finalizar la revisión), descargar \acs{pdf}, ver
    anotaciones, ver rúbrica y ver nota desglosada. Por defecto, todos los permisos están
    activados excepto la descarga del \acs{pdf} y el acceso extendido al \acs{pdf}. Estos
    permisos se aplican de forma uniforme a todas las instancias del examen.

    \textbf{Salida:} respuesta con la configuración de permisos actualizada.

    % RF-6.11 ---------------------------------------------------------------
    \item \textbf{Definición de rúbrica por problema.} \label{rf:exams:rubric}
    El sistema debe permitir definir una rúbrica de \textit{checkmarks} con puntuación fija
    para cada problema.

    \textbf{Entrada:} petición POST al recurso de rúbrica del problema con la lista de
    criterios, cada uno con descripción textual y puntuación fija (positiva o negativa).

    \textbf{Proceso:} el sistema almacena los criterios de la rúbrica vinculados al
    problema. La suma de puntuaciones positivas debería coincidir con la puntuación máxima
    del problema; se permite flexibilidad para incluir bonificaciones y penalizaciones. Los
    criterios con puntuación negativa se interpretan como penalizaciones.

    \textbf{Salida:} respuesta con la rúbrica completa del problema y código
    \acs{http}~201.

    % RF-6.12 ---------------------------------------------------------------
    \item \textbf{Listado de exámenes de una asignatura.} \label{rf:exams:list}
    El sistema debe permitir consultar la lista de exámenes de una asignatura.

    \textbf{Entrada:} petición GET al recurso de exámenes de la asignatura.

    \textbf{Proceso:} el sistema recupera los exámenes de la asignatura adaptando la
    información al rol del usuario: los estudiantes ven solo los exámenes cuyas instancias
    están publicadas o en estado posterior; los profesores y el coordinador ven todos los
    exámenes.

    \textbf{Salida:} respuesta con la lista de exámenes, incluyendo nombre, número de
    modelos y número de instancias.

    % RF-6.13 ---------------------------------------------------------------
    \item \textbf{Consulta detallada de examen.} \label{rf:exams:detail}
    El sistema debe permitir consultar la información completa de un examen, incluyendo sus
    modelos, problemas, rúbricas y estado.

    \textbf{Entrada:} petición GET al recurso del examen.

    \textbf{Proceso:} el sistema recupera los datos completos del examen, incluyendo la
    lista de modelos con sus problemas, puntuaciones, zonas, rúbricas y un resumen del
    estado de las instancias. La información se adapta al rol del usuario.

    \textbf{Salida:} respuesta con los datos completos del examen.

    % RF-6.14 ---------------------------------------------------------------
    \item \textbf{Almacenamiento de \acs{pdf} en blanco en MinIO.}
    \label{rf:exams:blank_pdf_storage}
    El sistema debe almacenar los ficheros \acs{pdf} en blanco de los modelos de examen en
    el servicio de almacenamiento de objetos.

    \textbf{Entrada:} evento interno disparado por la carga de un \acs{pdf} en blanco
    (\cref{rf:exams:blank_pdf}).

    \textbf{Proceso:} el fichero se almacena en MinIO en un \dfn{bucket} específico para
    plantillas, con una clave que incluye el identificador de la organización, la
    asignatura, el examen y el modelo. Se almacenan los metadatos del fichero (tamaño, tipo
    MIME, fecha de carga) en la base de datos.

    \textbf{Salida:} confirmación del almacenamiento con la referencia interna al objeto en
    MinIO.

\end{functional}
